\documentclass[12pt,a4paper]{article}
\usepackage{times}
\usepackage{amsmath}
\usepackage{amssymb}
\usepackage{graphicx}
\usepackage{cite}
\usepackage{url}
\usepackage{hyperref}
\usepackage{geometry}
\usepackage{setspace}
\usepackage{indentfirst}
\usepackage{fancyhdr}
\usepackage[font={it,footnotesize}]{caption}
\usepackage{sectsty}
\usepackage{booktabs}

% Section title formatting
\sectionfont{\fontsize{12}{15}\selectfont\bfseries}
\subsectionfont{\fontsize{12}{15}\selectfont\bfseries}

% Page style
\pagestyle{fancy}
\fancyhf{}
\fancyfoot[C]{\thepage}
\renewcommand{\headrulewidth}{0pt}

% Margins and spacing
\geometry{left=2.54cm,right=2.54cm,top=2.54cm,bottom=2.54cm}
\renewcommand{\baselinestretch}{1.5}
\setlength{\parindent}{0.5cm}

% Title
\title{\fontsize{14}{17}\selectfont\textbf{Прогресс в глубоком обучении для классификации изображений дистанционного зондирования: Комплексный обзор}}
\author{Лян Цзинъюй\\
        Северо-Западный политехнический университет\\
        Сиань, Шэньси, Китай\\
        \texttt{3445078775@qq.com}}
\date{}

\begin{document}

\maketitle

\begin{abstract}
Глубокие нейронные сети коренным образом изменили способ анализа спутниковых и аэрофотоснимков. При правильном проектировании современные иерархические архитектуры достигают беспрецедентных уровней точности. Они превосходно справляются со сложными задачами, варьирующимися от интерпретации сцен до точной локализации объектов и семантической аннотации на уровне пикселей. Три критических фактора лежат в основе этих достижений: извлечение многомасштабных признаков, механизмы внимания, выделяющие значимые области, и трансферное обучение из крупномасштабных наборов данных.

Тем не менее, сохраняются значительные препятствия. Аннотированные обучающие выборки остаются дефицитными. Обучение и вывод накладывают существенные вычислительные нагрузки, усложняя бортовое развертывание. Производительность снижается при изменениях освещения, погоды, сезона или конфигураций датчиков.

Исследователи изучают несколько перспективных направлений. Многие теперь предварительно обучают модели на огромных коллекциях неаннотированных изображений, снижая зависимость от ручной аннотации. Растет интерес к слиянию оптических, радарных, высотных и гиперспектральных наблюдений. Тем временем другие разрабатывают компактные сети без ущерба для точности. Эти коллективные усилия способствуют расширению практических применений в данной области.
\end{abstract}

\textbf{Ключевые слова:} глубокое обучение, дистанционное зондирование, классификация изображений, сверточные нейронные сети, классификация сцен, обнаружение объектов

\section{Введение}
\label{sec:introduction}
Технология дистанционного зондирования претерпела колоссальные изменения за последнее десятилетие. Спутниковые группировки, воздушные платформы и наземные датчики теперь ежедневно передают огромные объемы изображений высокого разрешения. Эти наборы данных фиксируют разнообразную спектральную, пространственную и временную информацию, поддерживая такие приложения, как картирование землепользования, городское планирование, реагирование на стихийные бедствия и экологический мониторинг. Однако рост объема и сложности данных дистанционного зондирования быстро превзошел традиционные методы классификации.

Многие годы исследователи в значительной степени полагались на рукописные признаки. Эксперты разрабатывали дескрипторы на основе спектральных откликов, текстурных паттернов и геометрических форм. Хотя эти методы показывали вполне приемлемые результаты в простых сценариях, они испытывали трудности в гетерогенных средах или при переменных условиях визуализации. Разработка признаков требовала глубоких знаний в предметной области и значительного времени. Адаптация к новым датчикам или задачам оказалась медленной и трудоемкой.

Глубокое обучение полностью революционизировало эту ситуацию. Нейронные сети изучают представления непосредственно из необработанных пикселей, устраняя ручную разработку признаков. Сверточные нейронные сети (CNN) оказались особенно эффективными, поскольку они естественным образом захватывают пространственные иерархии и контекстуальные отношения, которые упускают традиционные методы.

Изображения дистанционного зондирования заметно отличаются от повседневных фотографий. Они часто содержат десять или более спектральных каналов, с пространственным разрешением, значительно различающимся между платформами. Объекты появляются в vastly разных масштабах в сложных пространственных расположениях. Аннотация остается как дорогой, так и трудоемкой. Эти уникальные характеристики способствовали созданию специализированных сетевых архитектур и стратегий обучения.

Данный обзор фокусируется на ключевых разработках с 2020 по 2023 год. Раздел 2 охватывает фундаментальные концепции. Раздел 3 описывает архитектурные инновации. Раздел 4 рассматривает три основные области применения: классификацию сцен, обнаружение объектов и семантическую сегментацию. Раздел 5 обсуждает текущие проблемы и будущие направления исследований, а Раздел 6 завершает статью.

\begin{table}[htbp]
\centering
\caption{Обзор применения глубокого обучения в дистанционном зондировании}
\label{tab:applications}
\begin{tabular}{@{}lll@{}}
\toprule
\textbf{Применение} & \textbf{Основная задача} & \textbf{Вызов} \\
\midrule
Классификация сцен & Маркировка на уровне изображения & Многомасштабные признаки \\
Обнаружение объектов & Локализация и распознавание & Обнаружение малых объектов \\
Семантическая сегментация & Классификация на уровне пикселя & Дисбаланс классов \\
\bottomrule
\end{tabular}
\end{table}

\section{Основы}
\label{sec:background}
\subsection{Основы сверточных нейронных сетей}
CNN были специально разработаны для структурированных в сетку данных, таких как изображения. Типичные архитектуры стекают сверточные слои, операции пулинга и полносвязные слои, перемежаемые нелинейными функциями активации.

Сверточные слои служат детекторами признаков. Каждый фильтр скользит по входу, реагируя на локальные паттерны, такие как края, углы или текстуры. Поскольку один и тот же фильтр применяется везде, количество параметров остается управляемым. Затем слои пулинга уменьшают пространственные размеры, сохраняя важные активации, обеспечивая инвариантность к переносу и снижая вычислительную нагрузку. Через последовательные слои сети строят все более абстрактные представления. Наконец, полносвязные слои генерируют оценки классификации или прогнозы ограничивающих рамок. Выпрямленное линейное устройство (ReLU) стало стандартной функцией активации благодаря устойчивости к исчезновению градиентов и более высокой скорости обучения.

\subsection{Обучение и оптимизация CNN}
Процесс обучения минимизирует расхождение между прогнозами и истинными метками. Обратное распространение вычисляет градиенты, а оптимизаторы, такие как стохастический градиентный спуск (SGD) или адаптивная оценка моментов (Adam), итеративно обновляют веса. Различные техники улучшают стабильность обучения и ускоряют сходимость.

Аугментация данных вводит случайные повороты, отражения и цветовые смещения, подвергая сети большему разнообразию. Методы регуляризации, такие как dropout и затухание весов, предотвращают переобучение. Пакетная нормализация стандартизирует активации внутри каждого мини-пакета, стабилизируя процесс обучения. Трансферное обучение оказывается особенно ценным в дистанционном зондировании, где предварительно обученные веса ImageNet инициализируют сети перед тонкой настройкой на меньших, специфичных для домена наборах данных.

\subsection{Специфические вызовы дистанционного зондирования}
Изображения часто содержат более десяти спектральных каналов, требуя модификаций входных каналов. Расстояние выборки на поверхности земли значительно варьируется, что приводит к огромным различиям в размерах объектов. Семы семантически плотные, содержащие сложные комбинации дорог, зданий, растительности, транспортных средств и водных объектов. Наборы данных tendem to be smaller than those in general computer vision due to the high cost of producing precise pixel-level or object-level labels. Дисбаланс классов также является распространенным явлением; например, лесной покров может значительно превышать площадь аэропортов.

Эти характеристики означают, что стандартные модели, обученные на ImageNet, часто показывают низкую производительность, мотивируя исследования в области специфических для домена решений.

\section{Прогресс в архитектурах CNN для дистанционного зондирования}
\label{sec:developments}
\subsection{Извлечение многомасштабных признаков}
Объекты на аэро- и спутниковых изображениях варьируются от отдельных автомобилей до целых городов. Обычные сети с фиксированными рецептивными полями часто упускают мелкие детали или широкий контекст. Было разработано несколько эффективных стратегий для решения этой проблемы.

Пирамидальный пулинг применяет операции пулинга на нескольких масштабах, затем комбинирует результаты, одновременно захватывая локальную и глобальную информацию. Расширенная свертка использует разные скорости дилатации для достижения аналогичных целей без потери пространственного разрешения. Сети признаковых пирамид постепенно сливают семантически сильные признаки из глубоких слоев с признаками высокого разрешения из мелких слоев через боковые соединения, создавая многоуровневые представления. Результирующие признаки обладают семантической устойчивостью на всех масштабах, что помогает обнаружению малых объектов.

\subsection{Механизмы внимания}
Внимание стало повсеместным в современных архитектурах. Пространственное внимание учится фокусироваться на информационно богатых областях, улучшая дискриминативные признаки и подавляя фоновый шум. Канальное внимание, как продемонстрировано в Squeeze-and-Excitation, перераспределяет карты признаков в соответствии с глобальной важностью. Комбинация обоих в модулях, таких как модуль внимания сверточного блока, обычно дает наилучшие результаты.

В сложных сценах, таких как порты или аэропорты, внимание может игнорировать большие однородные области (вода, взлетно-посадочные полосы) и фокусироваться на элементах, определяющих категорию, таких как корабли или самолеты.

\subsection{Трансферное обучение и адаптация домена}
Учитывая дефицит аннотированных данных дистанционного зондирования, большинство исследователей начинают с предварительно обученных на ImageNet основ. Стандартная тонкая настройка приносит значительные выгоды, но разрыв домена между естественными фотографиями и видами сверху остается существенным. Адаптация домена с использованием состязательных методов, выравнивание корреляций и самообучающиеся методы с использованием псевдометок показывают перспективы в дальнейшем сокращении этого разрыва.

\subsection{Легкие сетевые архитектуры}
Бортовая обработка на спутниках или дронах требует сетей размером менее 100 МБ, способных завершить вывод в течение десятков миллисекунд. Обрезка, квантование, дистилляция знаний и эффективные варианты внимания помогают удовлетворить эти ограничения. Поиск нейронной архитектуры с учетом аппаратного обеспечения все чаще проектирует сети, оптимизированные для оборудования спутников или дронов.

\section{Вызовы и будущие направления}
\label{sec:challenges}
\subsection{Ограниченные обучающие данные}
Высококачественные аннотированные наборы данных остаются самым большим узким местом. Самообучающееся предварительное обучение на массивных коллекциях неаннотированных изображений привлекло значительное внимание. Фреймворки контрастного обучения, такие как MoCo, SwAV и DINO, вместе с подходами маскированного моделирования изображений, такими как MAE и SimMIM, создают представления, которые могут соответствовать или превосходить обучение с учителем на ImageNet при тонкой настройке для задач дистанционного зондирования.

Методы слабого и полу-обучения также показывают перспективы. Метки на уровне изображения, грубые наброски или точечные аннотации стоят значительно дешевле, чем полные маски пикселей, но все еще могут обучать эффективные модели сегментации при использовании осмотрительно.

\subsection{Вычислительная эффективность}
Бортовой вывод в реальном времени требует сетей размером менее 100 МБ, завершающих вывод в течение десятков миллисекунд. Обрезка, квантование, дистилляция знаний и эффективные механизмы внимания все вместе способствуют удовлетворению этих требований. Поиск нейронной архитектуры с учетом аппаратного обеспечения генерирует сети, настроенные для процессоров спутников и дронов.

\subsection{Слияние мультимодальных данных}
Оптические датчики не работают ночью или под облаками. Радиолокационный датчик с синтетической апертурой (SAR) и лидар (LiDAR) работают во всех погодных условиях, но лишены богатых спектральных деталей. Как эффективно комбинировать эти дополнительные источники остается открытым вопросом. Раннее слияние доказывает простоту, но уязвимо к сдвигу домена. Позднее слияние предлагает устойчивость, но упускает межмодальные взаимодействия. Промежуточное слияние с использованием перекрестного внимания или совместных архитектур трансформера кажется наиболее перспективным.

\subsection{Интерпретируемость}
Для практического развертывания заинтересованным сторонам необходимо понимать, почему сети классифицируют область как "поврежденную" после стихийного бедствия. Отображение активации класса, взвешенное по градиенту (Grad-CAM), и визуализация внимания стали стандартной практикой. Методы интерпретируемости на основе концепций, связывающие активации с понятными для человека концепциями, такими как индексы растительности или строительные материалы, появляются, но остаются незрелыми.

\section{Заключение}
\label{sec:conclusion}
С 2020 по 2023 год область достигла замечательного прогресса. Комбинация многомасштабного слияния признаков, механизмов внимания, улучшенного трансферного обучения и эффективных архитектур pushed производительность на уровни, позволяющие практическое развертывание. Однако дефицит данных, вычислительные ограничения, трудности слияния мультимодальных данных и пробелы в интерпретируемости все еще препятствуют более широкому внедрению.

Взгляд в будущее показывает, что самообучающиеся и фундаментальные модели, обученные на миллиардах неаннотированных спутниковых изображений, вероятно, сократят потребности в аннотации. Эффективные архитектуры, настроенные для граничных устройств, станут более распространенными. Интеллектуальное слияние оптических, SAR, гиперспектральных и временных данных откроет новые приложения. Методы интерпретируемости созреют, создавая доверие для высокорисковых решений.

Дистанционное зондирование генерирует экспоненциально растущие данные ежегодно. Глубокое обучение, адаптированное к уникальным вызовам этой области, остается самым мощным инструментом для преобразования массивных потоков пикселей в действенные инсайты.

\begin{table}[htbp]
\centering
\caption{Сравнение методов глубокого обучения в дистанционном зондировании}
\label{tab:comparison}
\begin{tabular}{@{}llll@{}}
\toprule
\textbf{Метод} & \textbf{Преимущества} & \textbf{Ограничения} & \textbf{Лучшие применения} \\
\midrule
Основанный на CNN & Высокая точность & Высокие требования к данным & Классификация сцен \\
Основанный на внимании & Выбор признаков & Высокие вычислительные затраты & Сложные сцены \\
Легкие модели & Высокая эффективность & Возможная потеря точности & Развертывание на границе \\
Мультимодальное слияние & Богатая информация & Сложная интеграция данных & Задачи слияния \\
\bottomrule
\end{tabular}
\end{table}

\bibliographystyle{IEEEtran}
\begin{thebibliography}{8}
\bibitem{zhu2017deep} Zhu, X. X., Tuia, D., Mou, L., et al. (2017). Deep learning in remote sensing: A comprehensive review and list of resources. \emph{IEEE Geoscience and Remote Sensing Magazine}, 5(4), 8-36.
\bibitem{cheng2016survey} Cheng, G., and Han, J. (2016). A survey on object detection in optical remote sensing images. \emph{ISPRS Journal of Photogrammetry and Remote Sensing}, 117, 11-28.
\bibitem{xia2017aid} Xia, G. S., Hu, J., Hu, F., Shi, B., Bai, X., Zhong, Y., and Lu, X. (2017). AID: A benchmark data set for performance evaluation of aerial scene classification. \emph{IEEE Transactions on Geoscience and Remote Sensing}, 55(7), 3965-3981.
\bibitem{maggiori2017end} Maggiori, E., Tarabalka, Y., Charpiat, G., and Alliez, P. (2017). Convolutional neural networks for large-scale remote-sensing image classification. \emph{IEEE Transactions on Geoscience and Remote Sensing}, 55(2), 645-657.
\bibitem{zhang2016weakly} Zhang, F., Du, B., Zhang, L., and Xu, M. (2016). Weakly supervised learning based on coupled convolutional neural networks for aircraft detection. \emph{IEEE Transactions on Geoscience and Remote Sensing}, 54(9), 5553-5563.
\bibitem{liu2019learning} Liu, Y., Chen, X., Wang, H., Ye, Z., Lin, Z., and Yu, P. S. (2019). Deep learning for pixel-level remote sensing data analysis: A review. \emph{ISPRS Journal of Photogrammetry and Remote Sensing}, 150, 1-15.
\bibitem{zeng2018improving} Zeng, D., Chen, S., Chen, B., and Li, S. (2018). Improving remote sensing scene classification by integrating global-context and local-object features. \emph{Remote Sensing}, 10(7), 734.
\bibitem{deng2017enhanced} Deng, Z., Sun, H., Zhou, S., Zhao, L., Lei, H., and Zou, H. (2017). Multi-scale object detection in remote sensing imagery with convolutional neural networks. \emph{ISPRS Journal of Photogrammetry and Remote Sensing}, 130, 83-94.
\end{thebibliography}
\end{document}